\documentclass{beamer}
\usepackage{beamerthemeCambridgeUS}
\usepackage{color}  % for \textcolor{red}{blah}
\usepackage{graphicx}
\usepackage{subfigure}
\usepackage{multicol}
\usepackage{verbatim} % for \begin{comment}
\usepackage{lipsum}


\newcommand\Fontvi{\fontsize{6}{7.2}\selectfont}


\newenvironment{packed_enum}{
\begin{enumerate}
  \setlength{\itemsep}{1pt}
  \setlength{\parskip}{0pt}
  \setlength{\parsep}{0pt}
}{\end{enumerate}}

\newenvironment{packed_item}{
\begin{itemize}
  \setlength{\itemsep}{1pt}
  \setlength{\parskip}{0pt}
  \setlength{\parsep}{0pt}
}{\end{itemize}}

\usetheme{CambridgeUS}

\title{Discrete Mathematics}
\author[Geoffrey Buhl]{Geoffrey Buhl \\ \texttt{geoffrey.buhl@csuci.edu}}
\institute[CSUCI]{California State University Channel Islands}
\date{week 4 class 2}

\def \vec{{\rm vec}}
\newtheorem{conjecture}[theorem]{Conjecture}
\newtheorem{explanation}[theorem]{Explanation}


\begin{document}

%\frame{\titlepage}

%\section[Outline]{}
%\frame{\tableofcontents}

\frame
{
  \frametitle{Stars and Bars Method}

The stars and bars method is a technique to count the number of ways to make $k$ choices from $n$ objects where order does not matter and repetition is allowed.
\vfill
Class Session Goals:
\begin{packed_enum}
\item Understand what it means to chose $k$ times from $n$ objects with repetition allowed and where order does not matter.
\item Present the Stars and Bars technique for counting
\item Present how to modify Stars and Bars technique
\item Practice Modifying the Stars and Bars technique
\end{packed_enum}

}

\frame
{
  \frametitle{Order and Repetition}

Consider making $k$ choices from a list of $n$ possible choices.  How many ways are there to do this?
\vfill

In order to be able to answer this question we need to establish some rules, and we do this by answering two questions:\\
\vspace{.25cm}
1. Does the order of our choices matter?\\
\vspace{.25cm}
2. Are choices allowed to be repeated?
\vfill

}


\frame
{
  \frametitle{Order and Repetition}

Order matters and repetition is allowed: $n^k$ \vfill

Order matters and repetition is not allowed: $n(n-1)(n-2)\cdots(n-k+1)=\frac{n!}{(n-k)!}$\vfill

Order does not matter and repetition is not allowed: $\binom{n}{k}=\frac{n!}{(n-k)!k!}$\vfill

Order does not matter and repetition is allowed:  Figure it out today!

}


\section{Class 7 - Unordered selections with repetition}

\frame
{
  \frametitle{Starter Activity}

Choose some partners around you to work on one of the following questions.\vfill

{\bf  Question A}: Suppose you have 3 flavors of ice cream.  How many different bowls of 2 scoops of ice cream are possible?  How does the answer change if you add a fourth flavor?
\vfill

{\bf Question B}: Suppose you have 3 boxes label 1, 2, 3 and two black balls.  How many different ways are there to place the balls in the boxes?  How does the answer change if you add a fourth box labeled 4?
\vfill

{\bf Question C}: How many solutions to the equation $x+y+z=2$ are there if $x$, $y$, and $z$  must be non-negative integers? For example the ordered triple $(x, y, z)=(1,0,1)$ is one solution. How many solutions are there to $w+x+y+z=2$ if $w$, $x$, $y$, and $z$  must be non-negative integers? 

}


\frame
{
  \frametitle{Stars and Bars Method}

\begin{figure}[htp]
  \begin{center}
    \subfigure{\includegraphics[scale=.45]{s&b.png}}
  \end{center}
\end{figure}
\vfill
There is a two stage process to uniquely create stars and bars figures like this one.\vfill

{\bf Step 1} Choose $k$ locations for stars * from among  $k+n-1$ locations available.\vfill

{\bf Step 2} Choose $n-1$ locations for bars | from among  remaining locations available.\vfill

The ``magic'' here is changing the problem into one where repetition of choices is not allowed.


}

\frame
{
  \frametitle{Stars and Bars Activity}

In small groups, work on the ``Modifying Stars and Bars'' worksheet. The goals of this worksheet are to:
\vfill
\begin{packed_enum}
\item Practice modifying the stars and bars technique to new problems
\item Applying the stage based counting process of selecting positions for start and then bars.
\end{packed_enum}


}

\frame
{
  \frametitle{Stars and Bars Answers}


\vfill
\begin{packed_enum}
\item $\binom{15}{12}$ %How many ways are there to distribute 12 balls into six numbered boxes if boxes 1 and 2 must be empty?
\item $\binom{14}{9}$ %How many ways are there to distribute 12 balls into six numbered boxes if box 1 must contain at least 1 ball and box 2 must contain at least 2 balls?
\item $\binom{14}{12}$ %How many ways are there to distribute 12 balls into six numbered boxes if every even numbered box must be empty?
\item $\binom{14}{9}$ %How many ways are there to distribute 12 balls into six numbered boxes if every odd numbered box must have at least one ball?
\item $\binom{6}{2}\binom{11}{3}$ %How many ways are there to distribute 12 balls into six numbered boxes if exactly two of the boxes must be empty?
\end{packed_enum}


}


\frame
{
  \frametitle{Order and Repetition}
  How many ways are there to make $k$ choices from $n$ objects?

If order matters and repetition is allowed: $n^k$ \vfill

If order matters and repetition is not allowed: $n(n-1)(n-2)\cdots(n-k+1)=\frac{n!}{(n-k)!}$\vfill

If order does not matter and repetition is not allowed: $\binom{n}{k}=\frac{n!}{(n-k)!k!}$\vfill

If order does not matter and repetition is allowed:  $\binom{n-1+k}{k}$\vfill


}

\begin{comment}
\frame
{
  \frametitle{Challenging counting problem}

This problem draws from all of chapter 1. A candy merchant carries bubble gum balls in seven different colors.
\vfill
How many ways can you buy 6 bubble gum balls if each one must have a unique color?
\vfill
How many ways can you buy 6 bubble gum balls and distribute one each to 6 kids (Alice, Bob, Carol, Diane, Ernie, Fergie)?
\vfill
How many ways can you buy 6 bubble gum balls and distribute one each to 6 kids (Alice, Bob, Carol, Diane, Ernie,Fergie) so that each kid gets a different color?

}
\end{comment}



\end{document}

